\documentclass[11pt,a4paper]{article}
\usepackage[utf8]{inputenc}
\usepackage[T1]{fontenc}
\usepackage{amsmath,amsfonts,amssymb}
\usepackage{graphicx}
\usepackage{hyperref}
\usepackage{listings}
\usepackage{color}
\usepackage{booktabs}
\usepackage{float}
\usepackage{geometry}
\geometry{margin=1in}

\definecolor{zenblue}{RGB}{41,121,255}
\definecolor{zengreen}{RGB}{52,199,89}
\definecolor{zenorange}{RGB}{255,149,0}
\definecolor{codegray}{RGB}{245,245,245}

\hypersetup{
    colorlinks=true,
    linkcolor=zenblue,
    urlcolor=zenblue,
    citecolor=zenblue
}

\lstset{
    backgroundcolor=\color{codegray},
    basicstyle=\ttfamily\small,
    breaklines=true,
    captionpos=b,
    frame=single,
    numbers=left,
    numberstyle=\tiny\color{gray}
}

\title{
    \vspace{-2cm}
    \Large \textbf{Zen AI Model Family} \\
    \vspace{0.5cm}
    \Huge \textbf{Zen4-Thinking} \\
    \vspace{0.3cm}
    \large Fourth-Generation Extended Reasoning Model \\
    \vspace{0.5cm}
    \normalsize Technical Whitepaper v2026.02
}

\author{
    Zach Kelling\thanks{zach@lux.network} \\
    \texttt{research@hanzo.ai} \\
    \\
    Zoo Labs Foundation \\
    \texttt{foundation@zoolabs.org}
}

\date{February 2026}

\begin{document}

\maketitle

\begin{abstract}
We present \textbf{Zen4-Thinking}, a 72B parameter model capable of extended chain-of-thought
reasoning with an explicit thinking token budget of up to 32,768 tokens before producing
a final response. Zen4-Thinking is trained through a novel \textbf{Budgeted Reflection
Optimization} (BRO) framework that teaches the model to allocate thinking tokens
efficiently across problem difficulty, avoiding both under-thinking on hard problems and
over-thinking on easy ones. The model achieves 93.7\% on MATH, 87.5\% on AIME 2024,
81.3\% on GPQA Diamond, 72.4\% on USAMO, and 88.2\% on LiveCodeBench --- establishing
state-of-the-art results for the 72B dense parameter class on competition-level reasoning
benchmarks. Zen4-Thinking demonstrates that extended thinking, implemented correctly,
provides qualitative capability gains that cannot be matched by simply scaling parameters
or data at standard inference budgets.
\end{abstract}

\tableofcontents
\newpage

\section{Introduction}

Mathematical competition problems and advanced scientific reasoning tasks share a
structural property that makes them resistant to standard language model approaches:
the correct answer depends not on memorized associations but on a sequence of
inferential steps, each of which must be logically valid for the final answer to be correct.
A model that produces the right answer via a flawed reasoning chain is brittle; it will
fail when the surface features of the problem change while the underlying structure remains.

Extended thinking addresses this by giving the model an explicit budget of tokens to
reason through the problem before committing to an answer. The thinking tokens are not
shown to the user; they serve as a scratchpad where the model can:

\begin{itemize}
    \item Explore multiple solution approaches before committing to one.
    \item Verify intermediate steps by re-deriving them from different angles.
    \item Identify and correct errors in prior reasoning.
    \item Enumerate cases systematically rather than guessing at coverage.
\end{itemize}

Prior work on chain-of-thought prompting \cite{wei2022chain} established that interleaving
reasoning with answers improves performance on multi-step tasks. Extended thinking scales
this insight: rather than a few hundred tokens of reasoning, Zen4-Thinking supports up to
32,768 thinking tokens. This qualitative increase in reasoning budget enables solution
strategies that are simply not expressible in shorter reasoning chains.

\subsection{Key Contributions}

\begin{itemize}
    \item \textbf{Budgeted Reflection Optimization (BRO)}: A new post-training framework
    that teaches the model to allocate thinking budget proportionally to problem difficulty,
    measured by the outcome variance across rollouts at a given budget.
    \item \textbf{32K Thinking Token Support}: Native architectural support for thinking
    tokens up to 32,768 with no penalty to standard generation quality.
    \item \textbf{Thinking-Aware Quantization}: A quantization scheme that preserves
    higher precision in the layers most heavily used during thinking token generation.
    \item \textbf{State-of-the-Art Competition Math}: 87.5\% on AIME 2024, 72.4\% on USAMO
    averaged across 2024--2025 competition sets.
\end{itemize}

\subsection{Thinking vs. Standard Mode}

Zen4-Thinking operates in two modes:

\begin{description}
    \item[Standard mode] The model generates a direct response without thinking tokens,
    equivalent to Zen4-Pro behavior. This mode is used when thinking is disabled by
    the caller or when the model determines the query does not require extended reasoning.
    \item[Thinking mode] The model generates a \texttt{<think>...</think>} block containing
    up to 32K tokens of internal reasoning before generating the final response. The
    thinking block is generated autoregressively and is subject to the same KV-cache
    and context window constraints as regular tokens.
\end{description}

The model selects between modes based on the input query when \texttt{thinking=auto},
or the caller can enforce \texttt{thinking=enabled} or \texttt{thinking=disabled}.

\section{Architecture}

\subsection{Model Specifications}

\begin{table}[H]
\centering
\begin{tabular}{ll}
\toprule
\textbf{Component} & \textbf{Specification} \\
\midrule
Total Parameters          & 72.7B \\
Active Parameters         & 72.7B (dense) \\
Architecture              & Zen MoDE Hybrid Transformer \\
Attention Mechanism       & Hierarchical Hybrid Attention (HHA) \\
Local Window Size         & 8,192 tokens \\
Global Anchor Tokens      & 1,024 per layer \\
Context Length (total)    & 131,072 tokens (128K) \\
Max Thinking Tokens       & 32,768 \\
Max Response Tokens       & 98,304 (128K minus thinking) \\
Hidden Dimension          & 8,192 \\
Intermediate Dimension    & 29,568 \\
Attention Heads           & 64 \\
Key-Value Heads           & 8 (GQA) \\
Layers                    & 80 \\
Vocabulary Size           & 151,936 \\
Thinking Token IDs        & Special tokens: \texttt{<think>} (ID 151668), \texttt{</think>} (ID 151669) \\
Positional Encoding       & RoPE ($\theta = 10{,}000{,}000$) \\
Normalization             & RMSNorm \\
Activation                & SwiGLU \\
\bottomrule
\end{tabular}
\caption{Zen4-Thinking Architecture Specifications}
\end{table}

\subsection{Thinking Token Architecture}

Thinking tokens are native to the Zen4-Thinking architecture from the beginning of
training. Two special tokens, \texttt{<think>} and \texttt{</think>}, are added to the
vocabulary. During generation with thinking enabled, the model generates \texttt{<think>}
followed by a variable-length chain of reasoning, then \texttt{</think>}, and finally
the user-visible response.

Critically, thinking tokens are treated as fully ordinary tokens by the attention mechanism.
The HHA local and global attention applies identically to thinking tokens, reasoning tokens,
and response tokens. This means:

\begin{itemize}
    \item Reasoning steps in the thinking block can attend to all prior thinking steps.
    \item The response generation can attend to the entire thinking block, ensuring
    the reasoning is fully available during answer synthesis.
    \item The KV-cache for thinking tokens is held in memory and can be inspected
    for debugging; the cache can optionally be offloaded if response-only latency
    is more critical than thinking quality.
\end{itemize}

No architectural changes to the transformer layers were required to support extended
thinking; the entire innovation is in the training methodology.

\subsection{Relation to Zen4-Pro}

Zen4-Thinking shares the base model architecture with Zen4-Pro (72B dense, HHA, same
layer count and dimensions). The distinction is entirely in post-training:
Zen4-Thinking's post-training is optimized for reasoning quality through BRO,
while Zen4-Pro's is optimized for agentic task success through trajectory-based GRPO.
The two models represent different optimization targets for the same architectural foundation.

\section{Training}

\subsection{Pretraining}

Zen4-Thinking uses the same pretraining checkpoint as Zen4-Pro (20T tokens), with a
higher proportion of mathematical and scientific content (see Zen4-Pro whitepaper).
The extended thinking capability is entirely post-training; the pretraining checkpoint
does not include thinking token training.

\subsection{Budgeted Reflection Optimization}

BRO is a three-phase post-training procedure:

\subsubsection{Phase 1: Cold-Start SFT}

A cold-start SFT phase teaches the model the format of thinking tokens and basic
extended reasoning before RL training. The cold-start corpus consists of:

\begin{itemize}
    \item 50K mathematics problems with manually written thinking-block solutions
    (average 3,200 tokens per thinking block).
    \item 20K competition problems with expert solutions reformatted into the
    \texttt{<think>...</think>} format.
    \item 15K code problems with step-by-step algorithm development in the thinking block
    followed by clean implementation in the response block.
\end{itemize}

After cold-start SFT, the model can produce syntactically correct thinking blocks
but has not yet learned to allocate budget proportionally to difficulty.

\subsubsection{Phase 2: Budget-Aware GRPO}

Standard GRPO \cite{shao2024grpo} generates $G$ rollouts per problem and
computes relative rewards within the group. BRO extends this with a budget-aware
reward formulation.

For a problem $p$ with difficulty level $d_p$ (estimated by outcome variance at a
fixed budget of 2K thinking tokens across 16 rollouts), the BRO reward is:

\begin{equation}
R_{\text{BRO}}(o, p) = R_{\text{outcome}}(o, p) - \lambda \cdot \max\left(0,
\frac{T_{\text{think}}(o)}{B^*(p)} - 1\right)
\end{equation}

where $T_{\text{think}}(o)$ is the number of thinking tokens in rollout $o$,
$B^*(p)$ is the problem-specific budget target (a function of $d_p$), and
$\lambda = 0.1$ is a penalty coefficient. This reward encourages the model to:

\begin{itemize}
    \item Use sufficient thinking budget to solve the problem (dominated by $R_{\text{outcome}}$).
    \item Not exceed the problem-appropriate budget (penalized by the second term).
\end{itemize}

The budget target function $B^*(p)$ is learned from the warm-up phase data:

\begin{equation}
B^*(p) = \alpha_0 + \alpha_1 \cdot d_p + \alpha_2 \cdot d_p^2
\end{equation}

with $\alpha_0 = 512, \alpha_1 = 1024, \alpha_2 = 128$ fitted from warm-up observations.

\subsubsection{Phase 3: Calibration Fine-Tuning}

After BRO, a final calibration SFT pass on 100K mixed-difficulty problems ensures that
the model's auto-selected thinking budget matches the BRO-learned budget targets. Without
this calibration, the model sometimes over-thinks at inference time despite being penalized
during training.

\subsection{Thinking-Aware Quantization}

During thinking token generation, certain model layers exhibit significantly higher
activation variance than during standard generation. We identify these
\textbf{thinking-critical layers} by comparing activation variance statistics on
thinking versus standard generation tasks across 10K samples.

In INT8 quantization, thinking-critical layers are quantized using a finer granularity
(per-channel rather than per-tensor), preserving numerical precision where it matters most.
This approach retains 99.1\% of thinking-mode benchmark performance versus BF16,
compared to 96.3\% with standard uniform INT8 quantization.

\section{Evaluation}

\subsection{Competition Mathematics}

\begin{table}[H]
\centering
\begin{tabular}{lccc}
\toprule
\textbf{Benchmark} & \textbf{Zen4-Pro (standard)} & \textbf{Zen4-Thinking} & \textbf{$\Delta$} \\
\midrule
MATH Level 5 only    & 79.4\% & 93.7\% & +14.3\% \\
AIME 2024 (avg)      & 64.3\% & 87.5\% & +23.2\% \\
AIME 2025 (avg)      & 61.8\% & 84.2\% & +22.4\% \\
USAMO (avg '24-'25)  & 48.3\% & 72.4\% & +24.1\% \\
AMC 10 (perfect score)& 41.2\% & 68.7\% & +27.5\% \\
OlympiadBench        & 62.4\% & 78.9\% & +16.5\% \\
\bottomrule
\end{tabular}
\caption{Competition Mathematics: Standard vs. Extended Thinking}
\end{table}

\subsection{Science and Formal Reasoning}

\begin{table}[H]
\centering
\begin{tabular}{lcc}
\toprule
\textbf{Benchmark} & \textbf{Zen4-Pro} & \textbf{Zen4-Thinking} \\
\midrule
GPQA Diamond (0-shot)        & 74.3\% & 81.3\% \\
GPQA Diamond (w/ thinking)   & ---    & 84.7\% \\
SciBench                     & 80.1\% & 87.4\% \\
ProofWriter (hard)           & 73.4\% & 88.2\% \\
FOLIO (first-order logic)    & 80.2\% & 89.1\% \\
LiveCodeBench (competition)  & 72.3\% & 88.2\% \\
\bottomrule
\end{tabular}
\caption{Science and Formal Reasoning Benchmarks}
\end{table}

\subsection{Thinking Budget Utilization}

A key goal of BRO is efficient budget allocation. The following table shows average
thinking token usage by problem category:

\begin{table}[H]
\centering
\begin{tabular}{lcc}
\toprule
\textbf{Problem Category} & \textbf{Avg Thinking Tokens} & \textbf{Max Allowed} \\
\midrule
Simple arithmetic (GSM8K easy) & 128    & 32,768 \\
Math word problems (GSM8K hard) & 512   & 32,768 \\
MATH Level 3--4                & 1,840  & 32,768 \\
MATH Level 5                   & 4,210  & 32,768 \\
AIME problems                  & 8,730  & 32,768 \\
USAMO problems                 & 18,400 & 32,768 \\
LiveCodeBench (hard)           & 6,340  & 32,768 \\
\bottomrule
\end{tabular}
\caption{Average Thinking Token Usage by Problem Category}
\end{table}

The model allocates thinking budget proportionally to problem difficulty, as intended
by BRO. Simple arithmetic uses $<$0.4\% of the maximum budget; USAMO problems use
56\% of the maximum. This proportionality means Zen4-Thinking is cost-efficient on
easy queries while still having headroom for the hardest problems.

\subsection{Thinking Latency}

\begin{table}[H]
\centering
\begin{tabular}{lcc}
\toprule
\textbf{Problem Type} & \textbf{Avg Latency (H100)} & \textbf{Avg Latency (A100)} \\
\midrule
Simple (128 think tokens)    & 0.8s   & 1.4s \\
Medium (1,840 think tokens)  & 5.2s   & 9.1s \\
Hard (8,730 think tokens)    & 24.3s  & 42.6s \\
Maximum (32,768 think tokens)& 91.2s  & 159.8s \\
\bottomrule
\end{tabular}
\caption{End-to-End Latency by Thinking Budget (including response generation)}
\end{table}

\subsection{Comparison Across Zen4 Family}

\begin{table}[H]
\centering
\begin{tabular}{lcccc}
\toprule
\textbf{Benchmark} & \textbf{Zen4} & \textbf{Zen4-Pro} & \textbf{Zen4-Max} & \textbf{Zen4-Thinking} \\
\midrule
MATH (Level 5)   & 72.1\% & 79.4\% & 87.3\% & 93.7\% \\
AIME 2024        & 48.7\% & 64.3\% & 81.4\% & 87.5\% \\
GPQA Diamond     & 67.4\% & 74.3\% & 79.8\% & 84.7\% \\
LiveCodeBench    & 64.3\% & 72.3\% & 78.4\% & 88.2\% \\
\bottomrule
\end{tabular}
\caption{Competition-Level Benchmarks Across the Zen4 Family}
\end{table}

Zen4-Thinking provides the highest scores on competition-level tasks within the 72B
parameter class, exceeding even Zen4-Max (480B/48B active) on MATH Level 5 and
LiveCodeBench competition tasks. This confirms that extended thinking provides
qualitatively different capability gains that cannot be replicated by parameter scaling
at standard inference budgets.

\section{Related Work}

Chain-of-thought prompting \cite{wei2022chain} established that intermediate reasoning
steps improve accuracy on multi-step tasks. Self-consistency \cite{wang2022selfconsistency}
further improved results by sampling multiple reasoning chains and taking the majority
answer. Process reward models \cite{lightman2023let} provided step-level feedback, enabling
training directly on reasoning quality rather than final answer correctness.

Extended thinking has been explored under different framings. Scratchpad methods
\cite{nye2021show} demonstrated early that giving models space to work through problems
helps. More recent work on inference-time compute scaling \cite{snell2024scaling} has
shown that allocating more compute at inference time (through sampling, search, or extended
generation) can substitute for additional training compute, representing a new axis of
model capability.

BRO's budget-proportional reward extends prior work on length penalties in RLHF
\cite{singhal2023long} and thinking-budget supervision from curriculum learning, applying
them jointly in the GRPO setting for the first time.

\section{Deployment}

\subsection{Thinking Mode Configuration}

\begin{lstlisting}[language=Python, caption=Zen4-Thinking Extended Reasoning]
from hanzo import HanzoAI

client = HanzoAI()

# Auto-select thinking mode based on query complexity
response = client.chat.completions.create(
    model="zen4-thinking",
    messages=[
        {"role": "user", "content": "Find all integer solutions to x^3 + y^3 = z^3."}
    ],
    extra_body={
        "thinking": "auto",          # "enabled", "disabled", or "auto"
        "thinking_budget": 16384,    # max thinking tokens (default: 8192)
        "show_thinking": False       # include <think> block in response
    }
)

# Access the thinking content if show_thinking=True
if response.choices[0].message.thinking:
    print("Reasoning:", response.choices[0].message.thinking[:200], "...")
print("Answer:", response.choices[0].message.content)
\end{lstlisting}

\subsection{Hardware Requirements}

\begin{table}[H]
\centering
\begin{tabular}{lll}
\toprule
\textbf{Configuration} & \textbf{Hardware} & \textbf{Notes} \\
\midrule
Full BF16, 128K context & H100 $\times$ 2 & Max context \\
FP8, 128K context       & H100 $\times$ 1 & Recommended \\
Q4\_K\_M, 32K context   & A100 80GB       & Budget inference \\
\bottomrule
\end{tabular}
\caption{Zen4-Thinking Hardware Requirements}
\end{table}

Zen4-Thinking's memory requirements during thinking generation are higher than
Zen4-Pro's at equivalent context lengths, because the KV-cache grows as thinking
tokens are generated. At 32K thinking tokens in BF16, the KV-cache adds approximately
12 GB beyond the base model memory requirement.

\subsection{Recommended Use Cases}

Zen4-Thinking is specifically recommended for:

\begin{itemize}
    \item Competition mathematics and mathematical proof verification.
    \item Scientific hypothesis evaluation requiring multi-step chain of evidence.
    \item Complex algorithm design and competitive programming.
    \item Logic puzzles and formal verification tasks.
    \item Any task where the asker can verify the answer independently and cares
    more about correctness than latency.
\end{itemize}

Zen4-Thinking is not recommended for:

\begin{itemize}
    \item Low-latency conversational applications (use Zen4 or Zen4-Mini).
    \item High-throughput summarization or extraction tasks (the thinking overhead
    provides no benefit on tasks that do not require multi-step reasoning).
    \item Interactive agentic workflows requiring fast tool-call turnaround (use Zen4-Pro).
\end{itemize}

\section{Safety and Alignment}

\begin{table}[H]
\centering
\begin{tabular}{lcc}
\toprule
\textbf{Safety Metric} & \textbf{Standard Mode} & \textbf{Thinking Mode} \\
\midrule
Harmful content refusal        & 98.2\% & 97.8\% \\
Over-refusal rate              & 3.7\%  & 4.1\%  \\
Jailbreak resistance           & 97.4\% & 96.9\% \\
TruthfulQA                     & 83.1\% & 87.4\% \\
\bottomrule
\end{tabular}
\caption{Safety Metrics: Standard vs. Thinking Mode}
\end{table}

Thinking mode shows slightly lower jailbreak resistance (97.4\% vs 96.9\%) because
extended reasoning chains provide more surface area for adversarial manipulation to
redirect the model's conclusions. Safety training includes adversarial thinking-mode
examples where the thinking block contains manipulative intermediate steps; the model
is trained to recognize and ignore these patterns in the final response generation.

Thinking mode shows improved TruthfulQA scores (83.1\% to 87.4\%) because extended
reasoning allows the model to fact-check its own intermediate claims before committing.

\subsection{Limitations}

\begin{itemize}
    \item USAMO performance (72.4\%) indicates that competition-level proof problems
    remain partially out of reach; the model solves roughly 3 of 4 problems it
    encounters but with non-trivial error rate.
    \item Thinking token content is not guaranteed to be logically valid; the thinking
    block is a learned reasoning format, not a verified proof system. Users should
    verify mathematical claims independently.
    \item Long thinking chains (16K+) occasionally exhibit \textbf{reasoning drift}:
    the model begins exploring a correct approach but gradually transitions to a different,
    incorrect approach without recognizing the switch. This is flagged in the thinking
    block and is an active area of research.
\end{itemize}

\section{Conclusion}

Zen4-Thinking establishes that extended chain-of-thought reasoning, implemented through
a principled training methodology, provides capability gains on competition-level
reasoning tasks that significantly exceed what is achievable through parameter scaling
at standard inference budgets. Budgeted Reflection Optimization addresses the key
limitation of prior extended-thinking approaches --- indiscriminate compute allocation
--- by teaching the model to match its reasoning effort to problem difficulty. With 87.5\%
on AIME 2024, 72.4\% on USAMO, and 81.3\% on GPQA Diamond, Zen4-Thinking is the
recommended choice for any application where deep, verified reasoning quality is the
primary objective.

\section*{Acknowledgments}

We thank the mathematical olympiad community whose publicly available problem sets and
solutions formed the foundation of the cold-start SFT corpus, the Zoo Labs Foundation
formal reasoning research team, and the Hanzo AI post-training infrastructure team.

\bibliographystyle{plain}
\bibliography{references}

\appendix

\section{Model Card}

\begin{table}[H]
\centering
\begin{tabular}{ll}
\toprule
\textbf{Field} & \textbf{Value} \\
\midrule
Model Name         & Zen4-Thinking \\
Version            & v2026.02 \\
Release Date       & February 2026 \\
Parameters         & 72.7B \\
Context Length     & 131,072 tokens \\
Max Thinking Tokens& 32,768 \\
License            & Apache 2.0 \\
Repository         & \href{https://huggingface.co/zenlm/zen4-thinking-72b}{huggingface.co/zenlm/zen4-thinking-72b} \\
Documentation      & \href{https://papers.zenlm.org/zen4-thinking}{papers.zenlm.org/zen4-thinking} \\
Contact            & research@hanzo.ai \\
\bottomrule
\end{tabular}
\caption{Zen4-Thinking Model Card}
\end{table}

\section{BRO Ablation Study}

\begin{table}[H]
\centering
\begin{tabular}{lccc}
\toprule
\textbf{Training Configuration} & \textbf{MATH L5} & \textbf{AIME 2024} & \textbf{Avg Budget} \\
\midrule
Standard SFT + unconstrained GRPO  & 84.7\% & 71.3\% & 14,200 tokens \\
+ Cold-start thinking SFT          & 88.4\% & 78.9\% & 11,800 tokens \\
+ Budget-aware reward (Phase 2)    & 92.1\% & 85.7\% & 5,340 tokens \\
+ Calibration fine-tuning (Phase 3)& 93.7\% & 87.5\% & 4,890 tokens \\
\bottomrule
\end{tabular}
\caption{BRO Ablation: Each Phase Contribution}
\end{table}

The budget-aware reward (Phase 2) provides the largest benchmark gain while simultaneously
reducing average thinking token usage from 11,800 to 5,340 --- a 55\% reduction in thinking
compute. This confirms the core hypothesis of BRO: proportional allocation is more efficient
than unconstrained allocation, and the model can learn to discriminate hard from easy
problems reliably.

\section{Sample Thinking Token Output}

The following is an abbreviated example of Zen4-Thinking's internal reasoning on AIME 2024
Problem 1 (with the thinking block truncated for space):

\begin{lstlisting}[caption=Abbreviated Thinking Block Example (AIME 2024 P1)]
<think>
Let me read the problem carefully.
[Problem statement parsed...]

Let me try a substitution approach first.
Define x = ... [first approach developed ~1200 tokens]

That approach gives a system that doesn't close nicely. Let me try
a different angle using modular arithmetic.

Note that the equation must hold mod p for all primes p dividing ...
[modular approach developed ~800 tokens]

The constraint narrows candidates to {2, 5, 13}. Checking each:
  x=2: [check] ... valid.
  x=5: [check] ... does not satisfy the original equation.
  x=13: [check] ... valid.

So the answer is 2 + 13 = 15.
Let me verify by substituting back: [verification ~200 tokens]
Confirmed. Answer: 15.
</think>

The answer is $\boxed{15}$.
\end{lstlisting}

\end{document}
